\documentclass[fleqn,numbers=noenddot,headinclude,%twoside,%1headlines,%
				11pt,a4paper,footinclude,%
				cleardoublepage=empty,abstractoff %
                ]{scrartcl}

\usepackage[hyphens]{url}
\usepackage{float}

\def\tomas{Tom\'{a}\v{s} Peterka}
\def\tomasid{peterto8}
\def\franta{Franti\v{s}ek Hejl}
\def\frantaid{hejl}
\def\jan{J\'{a}n Tka\v{c}\'{i}k}
\def\janid{tkacijan}
\def\frantaid{hejl}

\def\laboratorynr{5}
\def\laboratoryname{Ov\v{e}\v{r}en\'{i} metody CTEX}
\def\date{\today}


\input{preambl}

%\usepackage[subfigure]{tocloft}
%\renewcommand{\cftpartleader}{\cftdotfill{\cftdotsep}} % for parts
%\renewcommand{\cftchapleader}{\cftdotfill{\cftdotsep}} % for chapters
%\renewcommand{\cftsecleader}{\cftdotfill{\cftdotsep}} % for sections

\usepackage[sl,SF]{subfigure}
\usepackage{wrapfig}
\usepackage[hang,marginal]{footmisc}
\usepackage[justification=raggedright, margin=1cm,singlelinecheck=false,format=hang, labelfont={sl,small,rm},font={small,sf}, skip=10pt]{caption}
\renewcommand*\thesubfigure{\alph{subfigure}}

\DeclareMathOperator*{\argmax}{arg\,max}
\DeclareMathOperator*{\abs}{abs}

\makeatletter
\newcommand\wrapfill{\par
    \ifx\parshape\WF@fudgeparshape
    \nobreak
    \vskip-\baselineskip
    \vskip\c@WF@wrappedlines\baselineskip
    \allowbreak
    \WFclear
    \fi
}
\makeatother

\setlength{\intextsep}{6pt}

\usepackage{titlesec}
\titleformat{\section}{\large\bfseries}{\thesection}{1em}{}
\titleformat{\subsection}{\normalsize\bfseries}{\thesection}{1em}{}
\titleformat{\subsubsection}{\normalsize\bfseries}{\thesection}{1em}{}

\usepackage{amsmath}
\usepackage{graphicx}

\begin{document}

\makesavenoteenv[figure*]{figure}
\pagestyle{empty}
\onecolumn
%%Manuelle titelseite
\null  % Empty line
\nointerlineskip  % No skip for prev line
\vspace*{-1cm}
\begin{center}

\includegraphics[width=7cm,keepaspectratio=true]{./imgs/fit_logo.pdf}\\[0.3cm]
Czech Technical University\\
Faculty of Informatics

\vspace{1.2cm}

\Large MI-ROZ -- Rozpoznávání\\[0.5cm]
\laboratoryname \normalsize\\

\vspace{0.7cm}

\begin{tabular}{ccc}
\tomas & \franta & \jan \\
\tomasid @fit.cvut.cz & \frantaid @fit.cvut.cz & \janid @fit.cvut.cz \\
\end{tabular} \\[0.7cm]
\date
\end{center}
\vfill
\begin{abstract}
Práce popisuje implementaci algoritmu CTEX a jeho otestování na sérii zkušebních dat z ÚTIA (Ústav teorie informace a automatizace AV ČR).
\end{abstract}
\vfill
\tableofcontents

\clearpage

\pagestyle{scrheadings}
\setcounter{page}{1}

\section{Implementace CTex algoritmu}
\subsection{GB-FAB filter}
\label{sub:gb-fab}
lorem ipsum
% subsection  (end)

\subsection{SOM \& K-means shlukování}
\label{sub:som}
lorem ipsum
% subsection  (end)

\subsection{Multispace K-means}
\label{sub:multi-k-means}
lorem ipsum
% subsection  (end)

\subsection{Extrakce textur}
\label{sub:textury}
lorem ipsum
% subsection  (end)

\subsection{ASKM}
\label{sub:askm}
lorem ipsum
% subsection  (end)


\section{Vyhodnocení}
\label{sec:vyhodnoceni}
lorem ipsum
% section  (end)

\section{Závěr}
\label{sec:zaver}
lorem ipsum
% section  (end)
%\lstset{style=java}
%\begin{lstlisting}
%
%\end{lstlisting}

\end{document}
